% ===== PRÉAMBULE CANONIQUE — à coller VERBATIM en tête de CHAQUE .tex =====
\documentclass[11pt,a4paper]{article}
\usepackage[utf8]{inputenc}\usepackage[T1]{fontenc}\usepackage{lmodern}
\usepackage[french]{babel}
\usepackage{amsmath,amssymb,mathtools}\usepackage{icomma}
\usepackage{siunitx}\sisetup{locale=FR}  % \qty{}{}, \si{}, \ang{}
\usepackage{xcolor}\usepackage{colortbl}
\usepackage{tikz}\usepackage{pgfplots}\pgfplotsset{compat=1.18}
\usepackage[siunitx,europeanresistors]{circuitikz} % circuits (élec.)
\usepackage[most]{tcolorbox}\usepackage{enumitem}\usepackage{array,booktabs}
\usepackage[a4paper,margin=2.2cm]{geometry}
\usetikzlibrary{arrows.meta,calc,angles,quotes,positioning,patterns,%
  backgrounds,shapes.geometric,decorations.pathmorphing,decorations.markings,intersections}
\definecolor{primary}{RGB}{222,49,99}\definecolor{primarydark}{RGB}{150,22,55}
\definecolor{defcol}{RGB}{0,114,178}\definecolor{methcol}{RGB}{52,92,148}
\definecolor{excol}{RGB}{0,135,90}\definecolor{attncol}{RGB}{200,40,55}
\tcbset{boxbase/.style={enhanced,breakable,boxrule=0.9pt,arc=2.5pt,
  left=3mm,right=3mm,top=2.3mm,bottom=2.3mm,fonttitle=\bfseries,coltitle=white}}
% --- boîtes TUTORIEL (noms EXACTS, ne pas renommer) ---
\newtcolorbox{cle}[1][]{boxbase,colback=defcol!6,colframe=defcol,
  title={\textsc{À retenir}\ifx\relax#1\relax\relax\else\textmd{ : \itshape #1}\fi}}
\newtcolorbox{methode}[1][]{boxbase,colback=methcol!7,colframe=methcol,sharp corners,
  title={\textsc{Méthode}\ifx\relax#1\relax\relax\else\textmd{ --- \itshape #1}\fi}}
\newtcolorbox{exemplebox}[1][]{boxbase,colback=excol!5,colframe=excol,
  title={\textsc{Exemple}\ifx\relax#1\relax\relax\else\textmd{ : \itshape #1}\fi}}
\newtcolorbox{piege}[1][]{boxbase,colback=attncol!6,colframe=attncol,
  title={\textsc{Piège}\ifx\relax#1\relax\relax\else\textmd{ : \itshape #1}\fi}}
\newtcolorbox{remarque}[1][]{boxbase,colback=black!4,colframe=black!45,
  title={\textsc{Remarque}\ifx\relax#1\relax\relax\else\textmd{ : \itshape #1}\fi}}
% --- boîtes EXERCICE / CORRIGÉ (noms EXACTS, auto-comptées) ---
\newtcolorbox[auto counter]{exo}[1][]{boxbase,colback=primary!4,colframe=primary,
  title={\textsc{Exercice}~\thetcbcounter\ifx\relax#1\relax\relax\else\textmd{ --- \itshape #1}\fi}}
\newtcolorbox[auto counter]{corrige}[1][]{boxbase,colback=excol!4,colframe=excol,
  title={\textsc{Corrigé}~\thetcbcounter\ifx\relax#1\relax\relax\else\textmd{ --- \itshape #1}\fi}}
\newcommand{\vv}[1]{\vec{#1}}\newcommand{\ds}{\displaystyle}
\newcommand{\R}{\mathbb{R}}\newcommand{\N}{\mathbb{N}}\newcommand{\Z}{\mathbb{Z}}
% =========================================================================
\begin{document}

\begin{exo}[hiérarchie des couleurs dans le verre]
Dans le verre, les indices vérifient $n_{\text{bleu}}>n_{\text{vert}}>n_{\text{rouge}}$.
\begin{enumerate}[label=\alph*)]
  \item Classe les célérités $v_{\text{bleu}}$, $v_{\text{vert}}$, $v_{\text{rouge}}$ dans le verre.
  \item Lorsque la lumière blanche entre dans le verre, quelle couleur est la \textbf{plus déviée}
        (la plus proche de la normale) ?
\end{enumerate}
\end{exo}

\begin{exo}[une lumière rouge sur un prisme]
On envoie sur un prisme en verre une lumière \textbf{rouge monochromatique} (une seule longueur d'onde).
\begin{enumerate}[label=\alph*)]
  \item Cette lumière est-elle \emph{déviée} par le prisme ?
  \item Est-elle \emph{décomposée} en plusieurs couleurs ? Justifie en une phrase.
\end{enumerate}
\end{exo}

\begin{exo}[la profondeur apparente d'une pièce]
Une pièce de monnaie repose au fond d'un bassin, à la profondeur réelle $h=\qty{1.60}{m}$. L'eau a
pour indice $n_1=1{,}33$ ; on regarde presque à la verticale ($n_2=1$ pour l'air).

Calcule la profondeur apparente $h'$ à laquelle on voit la pièce.
\end{exo}

\begin{exo}[remonter à la profondeur réelle]
Un poisson semble nager à une profondeur apparente $h'=\qty{0.60}{m}$ pour un observateur regardant
presque à la verticale (eau $n_1=1{,}33$, air $n_2=1$).

À quelle profondeur réelle $h$ le poisson se trouve-t-il ?
\end{exo}

\begin{exo}[la latte brisée]
On plonge une latte (une règle) dans l'eau ; sa partie immergée semble « brisée » à la surface.
\begin{center}
\begin{tikzpicture}[scale=0.9,>=Stealth]
  \draw[black!55,line width=0.8pt] (-2.6,1.0)--(-2.6,-2.1)--(2.8,-2.1)--(2.8,1.0);
  \fill[defcol!22] (-2.6,0) rectangle (2.8,-2.1);
  \draw[primary,line width=1.2pt] (-2.6,0)--(2.8,0);
  \node[font=\scriptsize] at (2.4,0.45){air};
  \node[font=\scriptsize] at (1.8,-1.8){eau};
  \draw[black,line width=2pt] (-2.2,0.9)--(0.2,-1.7);
  \node[below,font=\scriptsize] at (0.2,-1.7){$A$ (réel)};
  \draw[black!45,dash pattern=on 2pt off 2pt] (0.2,0)--(0.2,-1.0);
  \fill[black!45] (0.2,-1.0) circle (1.3pt);
  \node[right,font=\scriptsize,black!55] at (0.28,-0.9){$A'$ (vu)};
\end{tikzpicture}
\end{center}
\begin{enumerate}[label=\alph*)]
  \item L'extrémité immergée paraît-elle \emph{plus proche} ou \emph{plus loin} de la surface qu'en
        réalité ?
  \item De façon générale, un objet immergé paraît-il plus \emph{gros} ou plus \emph{petit} ?
\end{enumerate}
\end{exo}

\end{document}
